\section{Experimental Setup}\label{sec:setup}

\paragraph{Hardware.}
Every controlled measurement ran on RunPod pods with two NVIDIA RTX PRO 6000 Blackwell Server Edition GPUs (sm\_120,
\dv{expl.gpu_total_mib}\,MiB reported per card), connected to the host over PCIe, at \dv{price.2xpro6000.usd_h}\,USD
per pod-hour (\dv{price.2xpro6000.eur_h}\,EUR). The exploratory pod had a cgroup memory limit of
\dv{expl.cgroup_max_gib}\,GiB. The designed reference (E9, not run) would use two B300 GPUs at
\dv{price.2xb300.usd_h}\,USD per hour, twice the single-GPU quote of \dv{price.1xb300.usd_h}\,USD. A single 80\,GB
A100 (\dv{price.1xa100.usd_h}\,USD per hour) is out of reach for two reasons: according to the exploratory notes the
pinned vLLM has no Ampere path for this model's sparse attention (Ampere support is an open feature request,
\datafile{audit/evidence/gh_issue_57144.json}), and neither checkpoint fits (Section~\ref{sec:system}). Every run
records the GPU model, driver, PCIe link generation and width, NUMA topology, CPU and cgroup limits
(\datafile{experiments/lib.sh}), and that record is what later explained part of the difference between hosts.

\paragraph{Software.}
All inputs are pinned in \datafile{stack/pins.env}: the vLLM nightly image by digest (vLLM \dv{expl.vllm_version}),
vllm-exl3 \texttt{d3cfd394}, ExLlamaV3 \texttt{5be88657}, the sfxnz recipe \texttt{460ac89c} and Diffbot's recipe at Hub
revision \texttt{a1cf1bd2}, verified file by file against \datafile{stack/recipe.lock.tsv}. Each variant is built
inside the pinned image by \datafile{stack/build/bootstrap_pod.sh} (or the Dockerfile wrapping it), which records
versions, git revisions, hashes of patched files and the full package list. Every variant shares the same serving
settings: tensor parallelism across the two GPUs, a context limit of 262,144 tokens, an fp8 KV cache of 4.25\,GiB, up to
64 sequences, 16,384 batched tokens per step, CUDA graphs for decode with capture sizes up to 256, prefix caching on,
and text only (\datafile{stack/serve.sh}).

\paragraph{Models.}
The EXL3 pack is \texttt{diffbot/DeepSeek-V4.1-Flash-EXL3-3bpw-2x-RTX-PRO-6000} at revision \texttt{a1cf1bd2}
(\dv{static.pack.total_gib}\,GiB), and the reference is \texttt{deepseek-ai/DeepSeek-V4.1-Flash} at \texttt{2cba9e42}
(\dv{static.orig.total_gib}\,GiB). Both are under the \dv{hub.orig.license} license as declared on the Hub
(\datafile{data/static/hub/}). Weights are downloaded on the pod at their pinned revisions and never redistributed.

\paragraph{Three sessions, three hosts.}
The campaign ran in three sessions on pods of this type, each on a different host. The first, on
\dv{lean.session.date}, lasted \dv{lean.session.minutes} minutes and ran E2, E8, E10 and E3. The second, on
\dv{lean.session2.date}, lasted \dv{lean.session2.minutes} minutes and ran E12, the E3 baseline again and the zero-copy
rung. The third lasted \dv{lean.session3.minutes} minutes and ran E13. Table~\ref{tab:pod-hardware} gives each host's
fingerprint. The GPUs, driver and PCIe link are the same, but the CPU model and the NUMA placement of the two GPUs are
not. The third session's record notes the first host's public IP, but its CPU and GPU bus ids differ, so it is a
different machine. The quality reference (E11) ran through DeepSeek's official API instead of on our own hardware.
Appendix~\ref{app:plan} lists the status of every experiment and how each session went.

% Generated by hostoffload.analysis.pipeline. Do not edit.
\begin{table}[!hb]
\centering
\small
\begin{fittable}{tab:pod-hardware}
\begin{tabular}{lp{2.4cm}p{2.2cm}rp{1.1cm}llp{1.2cm}p{1.6cm}}
\toprule
session & run & CPU & threads & GPU NUMA & GPU link & driver & PCIe max & bus ids \\
\midrule
1 (pod 1) & E8-main/main & AMD EPYC 9534 64-Core & 224 & 1, 1 & NODE & 595.91.07 & Gen5 x16 & 81:00.0, E1:00.0 \\
2 (pod 2) & E12-long-context/long & AMD EPYC 9535 64-Core & 256 & 0, 1 & SYS & 595.91.07 & Gen5 x16 & 06:00.0, 86:00.0 \\
3 (pod 3) & E13-coding-agent/agent & AMD EPYC 9535 64-Core & 256 & 1, 1 & NODE & 595.91.07 & Gen5 x16 & 86:00.0, F4:00.0 \\
\bottomrule
\end{tabular}
\end{fittable}
\caption{Hardware fingerprint of the three pod sessions (all 2x RTX PRO 6000 Blackwell Server Edition, same datacenter), from one run folder per session: CPU model and threads visible to the container (\texttt{fingerprint.json}), NUMA node of each GPU and the path between GPU0 and GPU1 (\texttt{nvidia-smi topo -m}: NODE within one NUMA node, SYS across the inter-socket link), driver, maximum PCIe link and GPU bus ids (\texttt{nvidia-smi -q}).}
\label{tab:pod-hardware}
\end{table}

